\documentclass[twocolumn,aps,prl,superscriptaddress,nofootinbib]{revtex4-2}
\usepackage{amsmath,amssymb,graphicx,hyperref,bm}

\begin{document}

\title{Replication Report: Hall mass and transverse Noether spin currents\\ in noncollinear antiferromagnets (Wernert \emph{et al.} 2024)}
\author{Automated Replication Subagent (Hermes)}
\affiliation{textures-100 replication campaign}
\date{\today}

\begin{abstract}
We report a from-scratch replication of the central analytic results of
Wernert, Pradenas, Tchernyshyov, and Chen, \emph{Hall mass and transverse
Noether spin currents in noncollinear antiferromagnets}
(arXiv:2404.12898v2). The paper's headline is that a \emph{static twist} of
the noncollinear kagome antiferromagnetic (AFM) order produces a purely
\emph{transverse} Noether spin current---a Hall-like response of the spin
current. Using a symbolic (sympy) re-derivation of the continuum Noether
current directly from the paper's Lagrangian and $\Gamma$-tensor, we reproduce
the headline \emph{exactly}: $J^y = \pm(\sqrt3/8)\,J S^2\,(\partial_x\phi)\,n_y$
with $J^x\equiv 0$ and a sign set by the direct/inverse triangular order. We
further reproduce the dynamical d.c.\ Hall response sign structure of Fig.~2
and the polycrystal Hall-mass magnon-velocity splitting of Eq.~(13). Verdict:
\textbf{REPLICATED} (Coverage 8/10, Agreement 9/10).
\end{abstract}

\maketitle

\section{Background and headline claim}
Noncollinear AFMs of the $\mathrm{Mn_3X}$ family host coplanar $120^\circ$
spin order on a kagome lattice. Because global spin-rotation symmetry is fully
broken by the order parameter, whether a \emph{conserved} spin current can even
be defined is nontrivial. Wernert \emph{et al.} apply Noether's theorem to the
continuum $\sigma$-model of the kagome Heisenberg AFM and show that the
conserved spin current generically acquires a component \emph{transverse} to
the gradient of the order parameter---a ``Hall spin current.'' The cleanest
manifestation is a \emph{static twist} of the spin frame,
\begin{equation}
\partial_x n_\alpha = (\partial_x\phi)\, n_x\times n_\alpha,
\qquad \partial_y n_\alpha = 0,
\end{equation}
which the paper states yields the purely transverse current
\begin{equation}
J^y = \pm\frac{\sqrt3}{8}\,J S^2\,(\partial_x\phi)\,n_y ,
\label{eq:headline}
\end{equation}
the $+/-$ sign corresponding to direct/inverse triangular order.

\section{Model and method}
The continuum Lagrangian and gradient tensor are
\begin{align}
\mathcal{L} &= \frac{\rho}{4}\,\partial_t n_\alpha\!\cdot\!\partial_t n_\alpha
 - \frac14\,\Gamma_{ab}^{\alpha\beta}\,\partial_a n_\alpha\!\cdot\!\partial_b n_\beta,\\
\Gamma_{ab}^{\alpha\beta} &= \eta\,\frac{\sqrt3}{4} J S^2
 \big(\delta_{a\alpha}\delta_{b\beta}+\delta_{a\beta}\delta_{b\alpha}\big),
 \quad \alpha,\beta\in\{x,y\},
\end{align}
vanishing if $\alpha=z$ or $\beta=z$; $\eta=+1$ for direct triangular order and
$\eta=+1(-1)$ for $\alpha=\beta\,(\alpha\neq\beta)$ for inverse triangular
order. Noether's theorem gives the spin current
\begin{equation}
J^a = -\tfrac12\,\Gamma_{ab}^{\alpha\beta}\, n_\alpha\times\partial_b n_\beta .
\label{eq:noether}
\end{equation}

Because the paper's core results are analytic, the faithful replication is a
\emph{symbolic} re-derivation rather than a black-box number. We implemented
Eqs.~(3)--(5) in \texttt{sympy}, built $\Gamma$ for both order types, and
evaluated Eq.~\eqref{eq:noether} for the static twist, the dynamical d.c.\
response $\langle J_a^\alpha\rangle=\Gamma_{ab}^{\alpha\gamma}P_b^\gamma$, and
the polycrystal isotropic tensor
$\bar\Gamma_{ab}^{\alpha\beta}=g_H(\delta_{a\alpha}\delta_{b\beta}
+\delta_{a\beta}\delta_{b\alpha})+g_0\delta_{ab}\delta_{\alpha\beta}$.
The itinerant-electron Kubo/Berry machinery of the shared
\texttt{gobel2024\_sd\_skyrmion\_kubo\_Lz\_kernel.py} (G\"obel 2024,
arXiv:2410.00820) is credited as the methodological sibling for
texture-induced Hall responses; the magnonic Noether current here is derived
independently.

\section{Results}
\paragraph{T3 --- Static twist (headline).}
Symbolic evaluation of Eq.~\eqref{eq:noether} under the static twist yields
$J^x=\mathbf{0}$ and $J^y = (\sqrt3/8)\,J S^2\,(\partial_x\phi)\,\hat e_y$ for
direct order and the negative for inverse order, matching
Eq.~\eqref{eq:headline} to symbolic zero, including that the current is
\emph{purely transverse}.

\paragraph{T4 --- Dynamical d.c.\ Hall response.}
$\langle J_y^y\rangle=\Gamma_{yx}^{yx}P_x^x=\pm(\sqrt3/4)J S^2 P_x^x$, with the
sign flipping between direct ($\mathrm{Mn_3Ir}$) and inverse ($\mathrm{Mn_3Sn}$)
order---reproducing the opposite-sign transverse currents of Fig.~2(a) vs (b).

\paragraph{T5 --- Polycrystal Hall mass.}
Diagonalizing the $3\times3$ dynamical matrix built from $\bar\Gamma$ gives one
longitudinal and two degenerate transverse magnon branches whose squared-velocity
splitting equals $2g_H/\rho$ exactly, confirming that the isotropic Hall mass
$g_H$ (Eq.~13, Fig.~1) controls the mode splitting and, being $O(3)$-invariant,
survives angular averaging---hence the Hall spin current is generic to
noncollinear AFMs and polycrystals. We do not reproduce the exact
velocity-branch \emph{labeling} of Eq.~(13) because our simplified elastic
matrix omits the LLG spin-rotation projection term; only the
convention-independent splitting magnitude is verified.

\section{Verdict and scoring}
\textbf{Verdict: REPLICATED.} The headline static-twist Hall spin current is
reproduced exactly, along with the dynamical sign structure and the Hall-mass
mode splitting.
\begin{itemize}
\item \textbf{Coverage: 8/10.} All three analytic pillars covered; the full
linearized-LLG FM/AFM bilayer strip numerics (Fig.~2 magnitudes, End Matter)
were not rebuilt.
\item \textbf{Agreement: 9/10.} T3, T4 exact; T5 splitting exact but the
Eq.~(13) branch labeling not matched.
\end{itemize}

\section*{Data and code}
Result JSON and the replication script are in \texttt{report/evidence/}
(\texttt{wernert2024\_result.json}, \texttt{wernert2024\_replication.py}).
Kubo/Berry sibling kernel credited: \texttt{gobel2024\_sd\_skyrmion\_kubo\_Lz\_kernel.py}.

\end{document}
